\section{Literature review}

\subsection*{Generative models for LOB and tick data}
Wiese et al.~\cite{wiese2020quant} introduced Quant GANs, a temporal
convolutional GAN for univariate return series with explicit attention
to long-tailed marginals. Subsequent work has explored W-GAN-GP
variants~\cite{takahashi2019modeling}, RNN-based generators for tick
series~\cite{koshiyama2020generative}, and conditional architectures
that take volatility regimes or macro indicators as conditioning input.
These methods all model price- or return-level data; we model
recurrence-plot tensors derived from the LOB state. Their architectures
could be substituted into our Stage~4 with minimal modification, and we
treat an ablation across generator backbones as a natural follow-on
rather than a competing contribution. To our knowledge, none of these
methods report a per-instrument training-stability scalar that is
comparable across markets, which is the empirical primitive we
introduce.

\subsection*{Recurrence plots in finance}
Recurrence plots (RPs) were introduced by Eckmann, Kamphorst and
Ruelle~\cite{eckmann1987} as a tool for visualizing the recurrence
structure of trajectories in the reconstructed phase space of a
dynamical system. Their quantification, recurrence quantification
analysis (RQA), was developed by Zbilut and Webber~\cite{zbilut1992}
and Marwan and colleagues~\cite{marwan2007}, who codified six standard
scalars over a binary RP matrix: recurrence rate ($RR$), determinism
($DET$), laminarity ($LAM$), longest diagonal line ($L_{\max}$),
trapping time ($TT$), and entropy of the diagonal-line-length
distribution ($ENTR$).

In financial econometrics, RPs have been applied predominantly to
univariate price or return series~\cite{strozzi2002recurrence}, with
the goal of detecting regime changes or non-stationarity. The closest
sibling to the present work is Safa et al.~\cite{safa2024}, who use
non-thresholded RPs (i.e., distance plots) computed from sliding
windows over L1 tick streams as the input to a $\beta$-VAE for
manipulation detection on NASDAQ equities. We differ from Safa et al.\
on three structural axes. First, we use full L2 depth (top-25 levels
per side) rather than best-bid/best-ask only, raising the state
dimension from $\mathbb{R}^2$ to $\mathbb{R}^{50}$ and turning the
recurrence plot into a fingerprint of depth-resolved liquidity geometry
rather than spread dynamics alone. Second, we bucket under three
parallel clocks (physical, event, and information time) and report
every downstream finding three times in parallel; their setup is
single-clock (physical). Third, our target is not anomaly detection on
a known label but the seed-to-seed reproducibility of a generative
model trained on the RPs themselves; this makes the RP a piece of
training data rather than a piece of feature engineering. In short,
Safa et al.\ treat the RP as a sensor; we treat it as a distribution.
